\chapter{Experiments}
\label{sec:experiments}

\section{Experimental Scope}

The goal of the experiments is to evaluate GammaZero as an implemented proof-search and data synthesis framework. The primary question is whether the system can run end-to-end: generate candidate actions, verify them with Lean, expand the subgoals introduced by skeleton actions, patch invalid generations, score partial proof signal, and finally produce Lean-verified proofs without placeholders when the search succeeds.

The main correctness metric is root solved rate without \texttt{sorry} or \texttt{admit}. However, root solved rate alone does not explain how the system behaves. Since GammaZero is designed to produce structured search trajectories for future policy optimization, the experiments also analyze the contribution of hierarchical search and the scoring signals extracted from generated candidates.

The evaluation therefore focuses on two aspects:
\begin{itemize}
    \item \textbf{Proof success:} whether the final reconstructed proof is accepted by Lean without placeholders.
    \item \textbf{Hierarchy and scoring signals:} how much skeleton expansion contributes beyond direct root solving, and how partial proof signals are scored.
\end{itemize}

These measurements are intended to describe the behavior of the implemented framework rather than to claim that every design choice is globally optimal. For example, comparing priority-queue search with a simpler expansion schedule can show whether the search order matters in this implementation, but it is not a universal theorem about best-first search.

The resulting search records are treated as scored proof-search trajectories. They are suitable for future policy optimization, but this thesis reports only the proof-success, ablation, scoring, and final-verification results obtained from the current implementation.

\section{Setup}

\paragraph{Lean environment.}
All experiments are conducted in Lean 4 using version \texttt{v4.29.0-rc1}. The project uses Mathlib \citep{mathlib2020}
from the \texttt{leanprover-community/mathlib4} repository, pinned to revision:
\begin{leanblock}
33a7291a345d718bca41f8ae8a9bae365d8f2a3e
\end{leanblock}
The Mathlib input revision is \texttt{v4.29.0-rc1}. This pinned environment ensures that all Lean
executions, diagnostics, and proof-state reports are evaluated under a fixed library version.

\paragraph{Hardware.}
Experiments are run on a Linux workstation with an NVIDIA RTX 3090 GPU with 24GB VRAM, 64GB
system RAM, and an Intel Xeon Silver 4114 CPU. The CPU has 20 logical threads, with 10 physical cores
and 2 threads per core. Lean execution is parallelized using multiple Lean workers.

\paragraph{Policy model.}
Candidate local proof actions and skeleton actions are generated using \texttt{Gemini 3 Flash Preview}. The model is used
as a fixed proposal mechanism for search-data generation; no additional fine-tuning or reinforcement
learning update is performed in this thesis. The model samples Lean code conditioned on serialized proof
states and task-specific prompts for local proof or skeleton action generation.

\paragraph{Search configuration.}
For each selected proof state, the search samples candidate actions subject to global and per-state
generation limits. Local proof attempts are tried before skeleton requests, and the priority queue controls
which open state is expanded next. The maximum search depth is set to \(D=10\), and the maximum
number of graph nodes per theorem is set to \(N=512\).

Lean execution uses \(4\) parallel workers. Each Lean call has a timeout of \(120\) seconds, and the
Lean heartbeat limit is set to \(100000\). The Sorrifier is allowed up to \(50\) patching cycles for a
failed generation.

\paragraph{Generation configuration.}
The maximum generation length is set to \(8192\) new tokens. Sampling uses temperature \(1.0\), with
the default nucleus-sampling setting of the backend API, approximately \(p=0.95\) when applicable.
The generated output may include a \texttt{<think>} block, but only the extracted Lean code block is
submitted to Lean for execution.

\section{Dataset}

The evaluation uses the validation and test splits of miniF2F-v2, introduced by
\citet{ospanov2025minif2fLeanRevisited} as a corrected version of the original miniF2F benchmark
\citep{zheng2021minif2f}. The original miniF2F benchmark contains 488 formalized olympiad-level
mathematics problems and has become a standard benchmark for neural theorem proving. However,
the revisited study identifies mismatches and errors between informal and formal statements in the
original benchmark and provides corrected versions with verified alignment between the informal
problem statements and the corresponding Lean formalizations.

In this thesis, we evaluate on the \texttt{miniF2F-v2s} validation and test splits, the simplified aligned version of miniF2F-v2. The dataset is suitable for
testing both local tactic generation and higher-level proof decomposition because the problems cover
competition-style mathematics such as algebra, number theory, and inequalities.

\section{Metrics}

We report metrics in two groups: proof success and scoring signals.

\paragraph{Proof success.}
The main correctness metric is root solved rate without placeholders. A root theorem is counted as solved only when the final reconstructed proof is accepted by Lean and contains no real \texttt{sorry} or \texttt{admit}. This metric is deliberately strict: placeholder-compilable proofs, partially solved skeletons, and patched scripts do not count as solved roots.

For external context, we compare against the strongest whole-proof prover reported by
\citet{ospanov2025minif2fLeanRevisited} on miniF2F-v2s. This comparison is not a strict
same-budget comparison, because their results use a different model, prover architecture, and
evaluation setup. It is
included as a real published reference point rather than as an internal ablation.

\paragraph{Scoring signals.}
GammaZero records syntactic and dependency-aware scores for generated proof fragments. These signals are used to describe partial progress during search-data generation, but final success is still counted only after Lean accepts the reconstructed proof without placeholders.

\section{Main Results}

Table~\ref{tab:minif2f_results} compares GammaZero with recent tree-search, whole-proof
generation, and Apollo-augmented provers on miniF2F-test. For GammaZero, the sample budget is
reported as a cumulative search-node budget rather than as independent proof samples. The token
budget is the average number of tokens consumed per solved instance under the corresponding
cumulative budget, with early stopping after successful verification.

\begin{table}[t]
\centering
\caption{
Comparison on miniF2F-test.
For GammaZero, the sample budget denotes the maximum search node budget.
The token budget is averaged over solved instances under each cumulative node budget.
}
\label{tab:minif2f_results}
\resizebox{\linewidth}{!}{
\begin{tabular}{@{}lcccc@{}}
\toprule
\textbf{Method} & \textbf{Model size} & \textbf{Sample budget} & \textbf{Token budget} & \textbf{miniF2F-test} \\
\midrule

\multicolumn{5}{@{}l}{\textit{Tree Search Methods}} \\
\midrule
Hypertree Proof Search \citep{lample2022hypertree} & 0.6B & $64 \times 5000$ & -- & 41.0\% \\
IntLM2.5-SP+BFS+CG \citep{wu2024internlm25stepprover} & 7B & $256 \times 32 \times 600$ & -- & 65.9\% \\
HPV16+BFS+DC \citep{hunyuanprover2025} & 7B & $600 \times 8 \times 400$ & -- & 68.4\% \\
BFS-Prover \citep{xin2025bfsprover} & 7B & $2048 \times 2 \times 600$ & -- & 70.8\% \\
\midrule

\multicolumn{5}{@{}l}{\textit{Whole-proof Generation Methods}} \\
\midrule
\makecell[l]{DeepSeek-R1-\\Distill-Qwen-7B \citep{guo2025deepseekr1}} & 7B & 32 & -- & 42.6\% \\
Leanabell-GD-RL \citep{leanabell2025prover} & 7B & 128 & -- & 61.1\% \\
STP \citep{yu2025stp} & 7B & 25600 & -- & 67.6\% \\
\midrule

\multirow{2}{*}{o3-mini \citep{openai2025o3mini}} & \multirow{2}{*}{--} & 1 & -- & 3.3\% \\
 &  & 32 & -- & 24.6\% \\
\midrule

o4-mini \citep{openai2025o3o4mini} & -- & 1 & -- & 7.0\% \\
\midrule

\multirow{3}{*}{Goedel-SFT \citep{lin2025goedelprover}} & \multirow{3}{*}{7B} & 32 & 16.1K & 57.6\% \\
 &  & 3200 & 1.6M & 62.7\% \\
 &  & 25600 & 12.7M & 64.7\% \\
\midrule

\multirow{3}{*}{\makecell[l]{Kimina-Prover-\\Preview-Distill \citep{wang2025kiminaprover}}} & \multirow{3}{*}{7B} & 1 & 4.4K & 52.5\% \\
 &  & 32 & 140K & 63.1\% \\
 &  & 1024 & 4.5M & 70.8\% \\
\midrule

\multicolumn{5}{@{}l}{\textit{Whole-proof Generation Methods + Apollo}} \\
\midrule
o3-mini + Apollo \citep{openai2025o3mini,apollo} & -- & 8 & -- & 40.2\% \\
o4-mini + Apollo \citep{openai2025o3o4mini,apollo} & -- & 15 & -- & 46.7\% \\
Goedel-SFT + Apollo \citep{lin2025goedelprover,apollo} & 7B & 362 & 179K & 65.6\% \\
Kimina-Preview + Apollo \citep{wang2025kiminaprover,apollo} & 7B & 307 & 1.3M & 75.0\%\\
\midrule
 
\multicolumn{5}{@{}l}{\textit{Our Method}} \\
\midrule
\multirow{5}{*}{\makecell[l]{GammaZero +\\Gemini 3 Flash Preview}}
& \multirow{5}{*}{--}
& 32 nodes
& 5.1K
& 59.8\% \\

&
& 64 nodes
& 8.4K
& 66.4\% \\

&
& 128 nodes
& 15.4K
& 73.8\% \\

&
& 256 nodes
& 19.8K
& 75.8\% \\

&
& 512 nodes
& 30.2K
& \textbf{77.5\%} \\

\bottomrule
\end{tabular}
}
\end{table}

GammaZero shows a clear budget-performance curve. With only 32 search nodes, it reaches 59.8\% on
miniF2F-test while using an average of 5.1K tokens per solved instance. Increasing the budget to 64
and 128 nodes raises the pass rate to 66.4\% and 73.8\%, respectively. The gains become smaller at
larger budgets: 256 nodes reaches 75.8\%, and 512 nodes reaches 77.5\%. This pattern suggests that
the hierarchical search finds many accessible proofs early, while the remaining problems require
substantially more search effort.

The comparison also highlights GammaZero's token efficiency. At 128 nodes, GammaZero obtains
73.8\% with 15.4K tokens per solved instance, exceeding Kimina-Prover-Preview-Distill at 1024
samples, which reports 70.8\% with 4.5M tokens. At 512 nodes, GammaZero reaches 77.5\%, improving
over Kimina-Preview + Apollo at 75.0\% while using 30.2K tokens per solved instance. Overall, the
main result supports the central claim that scaffold-aware hierarchical search can convert a small
cumulative node budget into competitive proof coverage.

All 59 solved roots pass final verification after proof stitching, so the reported
successes are placeholder-free Lean proofs. The search graph may contain patched candidates, temporary
placeholders, failed branches, and reserved skeletons, but a root is counted as solved only after the
selected proof path is reconstructed and checked again by Lean.

\section{Ablation Study}

To isolate the effect of hierarchical expansion, we compare GammaZero at the 32-node budget with the
full 512-node setting on the same evaluated test set. The 32-node setting corresponds to root-level
local proof search before deeper skeleton expansion becomes the main source of additional search
coverage. It is therefore an internal ablation, not an external theorem-proving baseline.

The root-only setting starts from the same initial theorem scaffold as GammaZero, with one root
\texttt{sorry} hole, and is reported as the cumulative 32-node budget in Table~\ref{tab:minif2f_results}.
The full setting keeps searching up to 512 nodes, allowing additional skeleton-created child states to
enter the proof graph. Both settings use the same verifier and final correctness criterion: a theorem is
counted as solved only when the reconstructed proof is accepted without \texttt{sorry} or \texttt{admit}.

\begin{table}[h]
\centering
\small
\begin{tabular}{@{}lccc@{}}
\toprule
Method & Solved & Solve Rate & Gain \\
\midrule
Root-only local proof (32 nodes) & -- & 59.8\% & -- \\
\textbf{GammaZero (512 nodes)} & -- & \textbf{77.5\%} & \textbf{+17.7\%} \\
\bottomrule
\end{tabular}
\caption{Root-only ablation. Gain is the absolute improvement from increasing the cumulative node budget from the root-only 32-node setting to the full 512-node hierarchical search setting.}
\label{tab:root-only-ablation}
\end{table}

The ablation shows that root-level local proof attempts solve a substantial fraction of the evaluated
theorems: the 32-node setting reaches 59.8\%. Increasing the budget to 512 nodes raises the pass rate
to 77.5\%, an absolute gain of 17.7 percentage points. This suggests that the hierarchical part of
GammaZero is not merely diagnostic machinery: it contributes additional solved roots beyond direct
completion at the initial proof state.

We further separate the solved theorems according to the depth of the final accepted proof path. A direct root solve has solved depth 0: the theorem is closed at the root node by
local proof actions without committing to a skeleton decomposition. A hierarchical solve has solved
depth at least 1: the final proof uses at least one skeleton-created child state in the committed
AND/OR proof tree.

\begin{table}[h]
\centering
\small
\begin{tabular}{@{}lcc@{}}
\toprule
Solve category & Budget & Rate \\
\midrule
Root-level local proof solves & 32 nodes & 59.8\% \\
Additional hierarchical search gain & 512 nodes & +17.7\% \\
\midrule
\textbf{Total solved} & \textbf{512 nodes} & \textbf{77.5\%} \\
\bottomrule
\end{tabular}
\caption{Budget-based decomposition of solved theorems on the evaluated test set. The 32-node setting corresponds to root-level local proof search; the additional gain is obtained by allowing deeper hierarchical search up to 512 nodes.}
\label{tab:test-depth-decomposition}
\end{table}

This budget split is the most direct internal measurement of the contribution of hierarchy. The
32-node result estimates the effective root-level baseline for the proposal model under the same
generation and verification conditions. The 17.7 percentage-point improvement at 512 nodes measures
the added coverage from enabling deeper skeleton decomposition and AND/OR graph search.

\begin{table}[h]
\centering
\scriptsize
\setlength{\tabcolsep}{4pt}
\begin{tabular}{@{}ll@{}}
\toprule
Theorem & Theorem \\
\midrule
\texttt{aime\_1983\_p9} & \texttt{amc12a\_2008\_p15} \\
\texttt{aime\_1988\_p3} & \texttt{amc12a\_2008\_p8} \\
\texttt{aime\_1990\_p2} & \texttt{amc12a\_2010\_p11} \\
\texttt{aime\_1996\_p5} & \texttt{amc12a\_2011\_p18} \\
\texttt{algebra\_amgm\_faxinrrp2msqrt2geq2mxm1div2x} & \texttt{amc12b\_2004\_p3} \\
\texttt{algebra\_amgm\_prod1toneq1\_sum1tongeqn} & \texttt{imo\_1961\_p1} \\
\texttt{amc12\_2000\_p15} & \texttt{imo\_1964\_p1\_1} \\
\texttt{amc12a\_2003\_p24} & \texttt{imo\_1964\_p1\_2} \\
 & \texttt{imo\_1965\_p1} \\
 & \texttt{imo\_1966\_p4} \\
\bottomrule
\end{tabular}
\caption{The 18 test theorems solved by hierarchical skeleton search. These theorems were not closed by direct root solving; their final accepted proofs required skeleton-created child states.}
\label{tab:hierarchical-test-solves}
\end{table}

\section{Patching and Scoring Signals}

Invalid Lean generations remain common in this setting. This is why binary success or failure is not
sufficient for data generation: many invalid candidates still contain locally meaningful structure that
can be patched and scored.

The low average dependency reward \(r_{\mathrm{dep}}\) should be read carefully. Most generated actions
do not produce a complete dependency structure: they may fail, be patched with placeholders, or leave
the main goal unsolved. Such actions naturally receive little or no dependency credit. A smaller number
of useful actions receive high dependency scores and are the actions that matter for successful stitched
proofs.

This is why \(r_{\mathrm{env}}\) and \(r_{\mathrm{dep}}\) are reported together. The syntactic reward measures
how much generated Lean text survives patching. The dependency reward is sparser: it gives credit only
when generated declarations are both solved and used by the checked proof. The average gap therefore
reflects the sparsity of genuinely useful proof structure, not a claim that every patched fragment is
semantically valuable.

A detailed trajectory case study for \texttt{aime\_1991\_p9} is included in Appendix~\ref{app:case-study-aime-1991-p9}.

\section{Discussion}

The main result shows that GammaZero can solve a large fraction of miniF2F-test with a compact
cumulative search budget. Its pass rate rises from 59.8\% at 32 nodes to 77.5\% at 512 nodes, while
the average token cost per solved instance remains far below the million-token regimes reported by
large-scale whole-proof sampling baselines. This supports the view that the graph can contain many
partial or failed branches while still producing a clean final proof when a complete path is found.

The patching and scoring results show that invalid generations are not rare. Most generated actions require some form of patching, but many still preserve enough local structure to receive useful scores. This supports the design decision to separate proof correctness from trajectory scoring. A patched script is not a proof, but it can still help describe how much of the generated Lean code was locally meaningful.

The root-only ablation shows that local proof generation is already strong, but not sufficient for the
evaluated subset. Increasing the budget from 32 to 512 nodes adds substantial coverage beyond direct
root completion, which supports the central design choice of combining local proof actions with
hierarchical search.

The gap between \(r_{\mathrm{env}}\) and \(r_{\mathrm{dep}}\) is one of the most important findings. It shows that local syntactic survival and final proof contribution are different properties. A generated fragment may survive patching but still not matter to the completed proof. This justifies using expression-tree dependency analysis as a separate structural signal.

At the same time, the results should not be overinterpreted. The external comparison uses published
systems with different models and budgets, while GammaZero is reported under the implementation
constraints of this thesis. The evaluation also
uses a fixed proposal model, a fixed Lean environment, and a finite search budget. Unsolved theorems
are not mathematically impossible; they are unsolved under the current configuration.


\section{Reproducibility Considerations}

Reproducibility is difficult in LLM-guided theorem proving because several components can change
independently. The model backend may update, sampling can be stochastic, Lean and Mathlib versions
can change theorem names or elaboration behavior, and worker timeouts can depend on hardware load.
For this reason, the experimental setup records the Lean version, Mathlib revision, model name, generation
parameters, worker count, timeout, search limits, and dataset split. These details are not administrative
noise; they define the environment in which a proof attempt is meaningful.

The most important reproducibility boundary is the Lean environment. A proof that verifies under one
Mathlib revision may fail under another because a lemma was renamed, a simp rule changed, or an
instance search path became different. The thesis therefore treats the pinned Lean and Mathlib versions
as part of the experimental artifact. Search records should be interpreted relative to that environment,
and future comparisons should either use the same environment or explicitly report the migration.

LLM sampling introduces a second source of variance. Even with fixed prompts and temperature, hosted
model APIs may not provide perfect determinism. The thesis therefore emphasizes aggregate search
metadata and trajectory structure in addition to root solve rate. If one run solves a theorem and another
does not, the difference may come from sampling rather than from the architecture. Recording generated
actions, priority scores, patching results, and final proof paths makes it possible to inspect why a run
succeeded or failed.

\section{Threats to Validity}

There are four main threats to validity. The first is benchmark validity. miniF2F-style datasets are useful
because they provide formal theorem statements, but they do not cover all styles of mathematics or all
Lean proof patterns. A system that works on olympiad algebra may behave differently on topology,
category theory, or large-library engineering proofs. This thesis therefore treats miniF2F-v2s as a
controlled evaluation setting rather than as a complete measure of theorem-proving ability.

The second threat is model dependence. GammaZero is designed to be model-agnostic, but the quality of
generated local proof actions and skeleton actions strongly affects the observed search behavior. A stronger model might
need fewer patches and produce better skeletons; a weaker model might make the search appear worse
than the architecture itself. The experiments should therefore report the proposal model clearly and avoid
claiming that results transfer automatically to all LLMs.

The third threat is heuristic dependence. Priority scores, skeleton commitment thresholds, and stale
skeleton rules are hand-designed in the current implementation. They are intended to be reasonable
engineering choices, not optimal search theory. The present experiments do not isolate each heuristic
component, so the exact numbers should be interpreted as evidence about this implementation rather than
universal constants. The root-only ablation isolates the broad contribution of hierarchical expansion,
but finer-grained ablations of the individual heuristic terms are left for future work.

The fourth threat is reward interpretation. Dense reward can make partial progress visible, but it can
also overemphasize syntactic survival if not paired with structural checks. GammaZero mitigates this by
combining line-based patching scores with dependency analysis and final Lean verification. Nevertheless,
the reward values should be understood as training signals, not as mathematical measures of elegance or
insight.
